\subsection{Fibonacci projection over a biquadratic character field}
\label{subsec:round4-m8-subfield}
The projection obstruction of
Proposition~\ref{prop:central-subcategory-projection} can be combined with
full-character induction without computing in an unnecessarily large
cyclotomic field.  We record the exact reduction used here.

\begin{proposition}[Projected twist traces]
\label{prop:round4-subfield-traces}
Let $K\subseteq L=\QQ(\zeta_F)$ be a Galois field containing all the
characters of a commutative fusion ring, and suppose that
$\operatorname{Gal}(K/\QQ)$ has exponent at most two.  Let $P$ be the
allowed prime support of the Frobenius--Schur exponent and $h\geq2$.
For a root of unity $\theta$, put
\[
 T_K(\theta)=\left(
 \frac{\operatorname{Tr}_{K(\theta)/K}(\theta^j)}{[K(\theta):K]}
 \right)_{1\leq j\leq h}.
\]
Every such vector with root order supported on $P$ is either zero or is
obtained from a root whose order divides
\[
 B=\prod_{p\in P}p^{\max(1,v_p(F))+\lfloor\log_p h\rfloor}.
\]
These vectors are computed exactly by first taking normalized trace to
$L$ and then averaging over $\operatorname{Gal}(L/K)$.  The induction
Gram, rotation, and low-character equations remain necessary after this
projection.  Their Galois multiplicity equalities may be imposed
separately for each projected twist vector.
\end{proposition}
\begin{proof}
The tower identity for normalized traces gives the stated computation.
The conductor bound follows from Lemma~\ref{lem:round3-relative-trace}:
if an order has a prime exponent larger than the displayed bound, every
power $1\leq j\leq h$ still has a nontrivial final cyclotomic layer over
$L$, and its relative trace vanishes.  Averaging once more preserves
zero.  Otherwise the order divides $B$.

All coefficients other than roots of unity in the induction and
low-character equations lie in $K$.  The normalized trace is
$K$-linear, so applying it preserves these equations.  The variables
counting central simple objects and cyclic eigenspaces remain
nonnegative; identifying equal projected vectors merely sums these
variables and thus weakens the necessary system.

For Galois balance, a modular Galois automorphism acts on joint character
eigenvalues by $\sigma$ and on the canonical twist by $\sigma^2$.
The latter automorphism restricts to the identity on $K$, by the
exponent-two hypothesis.  Normalized trace to $K$ is invariant under
this automorphism.  Consequently the projected twist vector is
unchanged along the corresponding permutation of central simples.
This proves the separate multiplicity equalities.  No integrality of
the projected trace, or of a solution to the relaxed linear system,
is asserted.
\end{proof}

For the ring $R_{6,8,3664}$, the Fibonacci basis object is $b_1$:
$b_1^2=1+b_1$.  Its Frobenius--Perron character field and global dimension
are
\[
 K=\QQ(\sqrt5,\sqrt6),\qquad
 D=300+132\sqrt5+120\sqrt6+54\sqrt{30}.
\]
The formal codegree polynomials are
\[
 t^2-15t+45,
 \qquad t^4-1200t^3+18000t^2-86400t+129600,
\]
each with multiplicity one.  Thus one may use $F=120$, $h=3$, and
$P=\{2,3,5\}$.  The bound is $B=720$, and there are exactly $43$
distinct projected vectors.  Arithmetic takes place in degree four,
whereas the ambient cyclotomic field has degree $32$.

By Corollary~\ref{cor:fibonacci-ising-projection}, each low column $c$
necessarily satisfies $c_1\leq c_0$.  For a high column, $c_0=0$;
if $c_1>0$, its twist must be $\zeta_5^{2}$ or $\zeta_5^{-2}$.
The two roots have the same normalized trace to the totally real field
$K$, so this condition leaves one projected twist vector for such a
column.  Columns with $c_1=0$ retain the full necessary trace domain.

\begin{theorem}
\label{thm:round4-3664}
The fusion ring $R_{6,8,3664}$ has no categorification over $\CC$.
\end{theorem}
\begin{proof}
Unrestricted signed sphericalization gives $768$ signed quotients and
$24$ sign orbits.  One orbit is gauge equivalent to the original ring;
the other orbits are eliminated by $36$ exact dimension-character
checks.  The unique surviving nonzero dimension-character component
is the Frobenius--Perron Galois orbit.  The general-to-pseudo-unitary
reduction therefore applies.

For the pseudo-unitary case, exhaustive low induction enumeration gives
$13$ frames.  Four fail an ordinary moment inequality; three fail the
Fibonacci twist-one projection inequality; four have no compatible
low-character labelling.  The remaining two frames each have one
labelling and one allowed second-indicator tuple.  Their necessary
systems have, respectively, $708$ equations in $2184$ variables and
$663$ equations in $1792$ variables.  Both systems have certified
rational separating vectors.  The verifier reconstructs every field,
frame, character assignment, indicator branch, and rational equation
from the full input tensor, and checks $A^Ty\geq0$, $b^Ty<0$ exactly.
The complete certificate is
\texttt{research/round4/m8/R6\_8\_3664\_certificate.json}, replayed by
\texttt{software/round4\_m8\_subfield.py}.
\end{proof}

The criterion was screened on all ten inputs to this continuation and
all $54$ known realizations.  The target is the only remaining input
containing a Fibonacci basis subcategory.  Eight realized rings contain
one; three lie in the implemented full-character subfield domain and
survive the necessary system, while the other five are explicitly
outside that domain.  The remaining $46$ controls have no Fibonacci
basis subcategory.  The structural census screen finds $443$ Fibonacci
subring cases among all $16612$ tensors; it records eligibility only,
not an additional $443$ exclusion replays.  Normal and optimized Python
replays of the new general certificate disable both numerical
optimizers.  Mutation checks separately reject omitted branches,
corrupted domain data, false projection witnesses, and altered
separating vectors.
